\documentclass[11pt]{article}
\usepackage[margin=0.7in]{geometry}
\usepackage{amsmath,amssymb}
\usepackage{enumitem}
\usepackage{tikz}
\usepackage{pgfplots}
\pgfplotsset{compat=1.18}
\usepackage{fancyhdr}
\usepackage{parskip}
\setlist[enumerate]{leftmargin=*, itemsep=6pt}
\pagestyle{fancy}
\fancyhf{}
\fancyfoot[L]{\small\textit{en}Vision\textsuperscript{\textregistered} Algebra 2}
\fancyfoot[C]{\small 1-4 Lesson Quiz, short form}
\fancyfoot[R]{\small\thepage}
\renewcommand{\headrulewidth}{0pt}
\newcommand{\choice}[1]{\tikz[baseline=-0.6ex]\node[draw, circle, inner sep=0.6pt, minimum size=3.3mm, font=\scriptsize\bfseries] {#1};\;}
\newcommand{\checkbox}{$\square$\;}
\newcommand{\answerbox}{\fbox{\rule{0pt}{14pt}\hspace{1.4cm}}}

\begin{document}

\noindent\textbf{Name}\; \rule{6cm}{0.4pt} \hfill Period \rule{1.2cm}{0.4pt} \quad Date \rule{2.2cm}{0.4pt}

\vspace{4pt}\noindent\rule{\textwidth}{0.6pt}

\vspace{2pt}\noindent{\LARGE\textbf{1-4 \; Lesson Quiz: Arithmetic Sequences and Series}} \hfill {\small short form}

\vspace{4pt}\noindent\rule{\textwidth}{0.6pt}

\vspace{6pt}
\textbf{1.} In November, the daily low temperature decreases by $2^{\circ}$F each day. Which of the following arithmetic sequences could represent the daily low temperature?

\vspace{4pt}
\begin{itemize}[leftmargin=1.2em, itemsep=5pt]
\item[] \choice{A} $2,\,-4,\,8,\,-16,\,32,\,\ldots$
\item[] \choice{B} $-2,\,-4,\,-8,\,-16,\,-32,\,\ldots$
\item[] \choice{C} $8,\,6,\,4,\,2,\,0,\,\ldots$
\item[] \choice{D} $-8,\,-6,\,-4,\,-2,\,0,\,\ldots$
\end{itemize}

\vspace{10pt}\noindent\rule{\textwidth}{0.4pt}

\vspace{6pt}
\textbf{3.} A bucket holds 100 potatoes. Hal removes and cleans 4 potatoes per minute. Which explicit definition matches the sequence that represents the potatoes left in the bucket after $n$ minutes?

\vspace{4pt}
\begin{minipage}[t]{0.40\textwidth}
\begin{itemize}[leftmargin=1.2em, itemsep=5pt]
\item[] \choice{A} $a_n=4-100(n-1)$
\item[] \choice{B} $a_n=100-4(n-1)$
\item[] \choice{C} $a_n=96+4(n-1)$
\item[] \choice{D} $a_n=96-4(n-1)$
\end{itemize}
\end{minipage}\hfill
\begin{minipage}[t]{0.56\textwidth}
Show your thinking:

\vspace{6pt}
$a_1=$ \rule{1.6cm}{0.4pt}

\vspace{8pt}
$d=$ \rule{1.6cm}{0.4pt}
\end{minipage}

\vspace{10pt}\noindent\rule{\textwidth}{0.4pt}

\vspace{6pt}
On Gabriela's first birthday, her parents gave her a \$50 savings account. On every birthday after that her parents added to the account, increasing the amount they deposited by \$25 each year.

\vspace{4pt}
\textbf{4.} Complete the recursive definition for the sequence that represents the amount of money that Gabriela's parents deposit into the account on her $n$th birthday.

\vspace{8pt}
\hspace{1.5cm}$f(n)=\begin{cases}\answerbox\,, & n=1\\[8pt] f(n-1)+\answerbox\,, & n>1\end{cases}$

\vspace{10pt}
What does the second number you wrote mean in this situation?

\vspace{4pt}\rule{\linewidth}{0.4pt}

\vspace{10pt}\rule{\linewidth}{0.4pt}

\vfill
\noindent\textbf{Scoring:} Question 3 and Question 4 are 10 points each. Question 1 is a bonus of up to 5 points. Total out of 20. (Questions 2 and 5 are on the full Savvas quiz only.)

\newpage
\noindent{\LARGE\textbf{1-4 Lesson Quiz, short form: answer key}}

\vspace{4pt}\noindent\rule{\textwidth}{0.6pt}

\begin{enumerate}
\item[\textbf{1.}] \textbf{C, $8, 6, 4, 2, 0, \ldots$} Each term is 2 less than the one before, so $d=-2$: the temperature decreases by $2^{\circ}$F each day. A and B are not arithmetic (there is no common difference: A multiplies by $-2$, B doubles); D is arithmetic but $d=+2$, so the temperature would increase. Savvas quiz item 1, DOK 3, skill F79. Prerequisite for Example 3 (the common difference; Example 1, cut). \textbf{Bonus 5:} 5 for C; 0 otherwise.
\item[\textbf{3.}] \textbf{D, $a_n=96-4(n-1)$.} After 1 minute $100-4=96$ potatoes are left, so $a_1=96$; 4 are removed each minute, so $d=-4$. Check: $a_2=96-4=92$ (after 2 minutes); $a_{25}=96-96=0$. Distractor B starts the count at 0 minutes ($a_1=100$); C has the wrong sign of $d$; A swaps $a_1$ and $d$. Savvas quiz item 3, DOK 2, skill F79. Taught by Example 3A (the explicit form $a_n=a_1+d(n-1)$) and the packet's Practice a (the same structure). \textbf{10 points:} 6 for D, 2 for $a_1=96$, 2 for $d=-4$.
\item[\textbf{4.}] \textbf{50; 25.} $f(n)=\begin{cases}50, & n=1\\ f(n-1)+25, & n>1\end{cases}$ The first deposit is \$50 and each deposit is \$25 more than the previous one, so the common difference is 25. Savvas quiz item 4, DOK 2, skill F79. Taught by Example 3B (the recursive form) and the packet's Practice b. \textbf{10 points:} 4 for 50, 4 for 25, 2 for a sentence saying each birthday's deposit is \$25 more than the year before (an explain prompt added to the TE item; it asks for no new computation).
\end{enumerate}

\textbf{Items left off this form} (available for a retake or Form B): Savvas item 2, Rayna's tapered staircase, total length of the stair tops from $3\tfrac12$ ft down to $2\tfrac16$ ft in steps of 4 in.\ (answer $\tfrac{85}{6}$ ft: 5 stairs, $S_5=\tfrac{5(\frac72+\frac{13}{6})}{2}=\tfrac{85}{6}$; DOK 2, F79; Examples 5-7, sums, cut); Savvas item 5, Gabriela's total deposits by her 18th birthday (answer \$4725: $a_{18}=50+25(17)=475$, $S_{18}=\tfrac{18(50+475)}{2}=4725$; DOK 2, F79; Example 7, cut). Both are arithmetic-series items; Topic 1 Assessment item 9 is the same skill.

\textbf{Realize PDF:} the TE items 3 and 4 are answer-only. If the Realize PDF is given instead of this form, score item 3 as 10 for D and item 4 as 5 for 50 and 5 for 25.

\textbf{IXL replacement:} SmartScore 70 or above on C8R, H82 and ZAH replaces this quiz with 20/20.

\end{document}
